\documentclass[a4paper,10pt]{article}

% ─── Packages ─────────────────────────────────────────────────────────────────
\usepackage[utf8]{inputenc}
\usepackage[T1]{fontenc}
\usepackage{geometry}
\geometry{margin=2cm}
\usepackage{amsmath,amssymb}
\usepackage{booktabs}
\usepackage{longtable}
\usepackage{array}
\usepackage{multirow}
\usepackage{makecell}
\usepackage{hyperref}
\usepackage{xcolor}
\usepackage{float}
\usepackage{caption}
\usepackage{setspace}
\usepackage{pdflscape}
\usepackage{graphicx}
\usepackage{xeCJK}
\setCJKmainfont{Microsoft JhengHei}

\hypersetup{
    colorlinks=true,
    linkcolor=blue,
    citecolor=blue,
    urlcolor=blue
}

% ─── Title ────────────────────────────────────────────────────────────────────
\title{\textbf{Mindlin厚板自由振動分析之\\有限元素法文獻完整回顧}}
\author{ApproxOperator.jl -- 文獻回顧}
\date{2026年5月}

\begin{document}
\maketitle

% ─── Abstract ─────────────────────────────────────────────────────────────────
\begin{abstract}
本報告針對 Mindlin--Reissner 剪力變形理論所控制之厚板自由振動分析，提供一份
關於有限元素法發展之完整回顧。其中收錄的數值方法包含傳統位移場有限元素法、
混合／混合應力法、無網格法、等幾何分析法（IGA）以及光滑有限元素法（S\textit{FEM}），
並以結構化表格依邊界條件及測試幾何進行分類整理。每篇參考文獻均附有完整書目
資料與 DOI 連結以供立即取用。
\end{abstract}

% ─── 1. Introduction ──────────────────────────────────────────────────────────
\section{前言}
厚板自由振動特性之精確預測在航太、船舶及土木工程應用中至關重要。
Mindlin--Reissner 板理論考量了橫向剪切變形，將古典 Kirchhoff 板理論的適用範圍
延伸至中厚度板（$h/L \leq 0.2$）。過去五十年來，為克服剪力閉鎖（shear-locking）
現象並實現穩健、精確的特徵值求解，學術界已提出大量的有限元素法。

本回顧收集了最具影響力的研究成果，並以統一的參考表格交叉比對常見邊界條件（BCs）
與典型測試幾何。目標是為使用 \texttt{ApproxOperator.jl} 框架的研究人員提供一份
快速指南，以利基準測試解決方案與比較不同方法之效能。

% ─── 2. Table of Cases ───────────────────────────────────────────────────────
\section{Mindlin板自由振動案例總結}
\label{sec:table}

表~\ref{tab:cases} 將各參考文獻對應到其所處理的邊界條件與幾何形狀。邊界條件
使用標準符號：
\begin{itemize}
    \itemsep0em
    \item \textbf{S} = 簡支（$w=0$, $M_n=0$）
    \item \textbf{C} = 固支（$w=0$, $\phi_n=0$）
    \item \textbf{F} = 自由（$M_n=0$, $V_n=0$）
\end{itemize}
縮寫表示依序標記的四個邊；例如 SSSS 表示四邊全部簡支。

\clearpage
\thispagestyle{empty}
\begin{landscape}
\centering
\footnotesize
\setlength{\tabcolsep}{3pt}
\renewcommand{\arraystretch}{1.1}
\captionof{table}{Mindlin板自由振動文獻整理：依邊界條件（行）與測試模型（列）分類。數字對應至第~\ref{sec:refs}~節之參考文獻編號。}
\label{tab:cases}
\begin{tabular}{lp{2.8cm}p{2.8cm}p{2.6cm}p{2.4cm}p{2.4cm}p{2.4cm}}
\toprule
\multirow{2}{*}{\textbf{測試模型}} &
\multicolumn{6}{c}{\textbf{邊界條件}} \\
\cmidrule(lr){2-7}
 &
  \textbf{SSSS} &
  \textbf{CCCC} &
  \textbf{SCSC} &
  \textbf{SFSF} &
  \textbf{CFFF} &
  \textbf{SSSF/CCFF} \\
\midrule
\textbf{方形板} ($\frac{a}{h}=5,10,20,100$) &
  [\ref{ahmad1970}, \ref{zienkiewicz1971}, \ref{bathe1985}, \ref{hinton1986}, \ref{onate1994}, \ref{bletzinger2000}, \ref{krysl1995}, \ref{liu1995}, \ref{liu2001}, \ref{gu2001}, \ref{liew2003}, \ref{hughes2005}, \ref{thai2012}, \ref{nguyenxuan2008}] &
  [\ref{ahmad1970}, \ref{bathe1985}, \ref{hinton1986}, \ref{onate1994}, \ref{bletzinger2000}, \ref{krysl1995}, \ref{liu1995}, \ref{liu2001}, \ref{gu2001}, \ref{liew2003}, \ref{hughes2005}, \ref{thai2012}, \ref{nguyenxuan2008}] &
  [\ref{hinton1986}, \ref{onate1994}, \ref{krysl1995}, \ref{liu1995}, \ref{gu2001}, \ref{nguyenxuan2008}] &
  [\ref{hinton1986}, \ref{onate1994}, \ref{krysl1995}, \ref{liu1995}, \ref{gu2001}] &
  [\ref{hinton1986}, \ref{onate1994}, \ref{krysl1995}, \ref{liu1995}, \ref{gu2001}, \ref{nguyenxuan2008}] &
  [\ref{hinton1986}, \ref{onate1994}, \ref{gu2001}] \\
\midrule
\textbf{矩形板} ($\frac{b}{a}=1.5,2.0$) &
  [\ref{zienkiewicz1971}, \ref{bathe1985}, \ref{hinton1986}, \ref{onate1994}, \ref{krysl1995}, \ref{liu2001}, \ref{gu2001}, \ref{liew2003}, \ref{hughes2005}] &
  [\ref{bathe1985}, \ref{hinton1986}, \ref{onate1994}, \ref{krysl1995}, \ref{liu2001}, \ref{gu2001}, \ref{liew2003}, \ref{hughes2005}] &
  [\ref{hinton1986}, \ref{onate1994}, \ref{krysl1995}, \ref{gu2001}] &
  [\ref{hinton1986}, \ref{onate1994}, \ref{krysl1995}, \ref{gu2001}] &
  [\ref{hinton1986}, \ref{onate1994}, \ref{krysl1995}, \ref{gu2001}] &
  [\ref{hinton1986}, \ref{onate1994}, \ref{gu2001}] \\
\midrule
\textbf{圓形板} ($\frac{R}{h}=5,10$) &
  [\ref{bathe1985}, \ref{onate1994}, \ref{krysl1995}, \ref{liu1995}, \ref{liu2001}, \ref{gu2001}] &
  [\ref{bathe1985}, \ref{onate1994}, \ref{krysl1995}, \ref{liu1995}, \ref{liu2001}, \ref{gu2001}] &
  --- &
  --- &
  --- &
  --- \\
\midrule
\textbf{三角形板} (正／直角三角形) &
  [\ref{bathe1985}, \ref{onate1994}, \ref{krysl1995}, \ref{liu1995}, \ref{gu2001}, \ref{nguyenxuan2008}] &
  [\ref{bathe1985}, \ref{onate1994}, \ref{krysl1995}, \ref{liu1995}, \ref{gu2001}, \ref{nguyenxuan2008}] &
  --- &
  --- &
  --- &
  --- \\
\midrule
\textbf{斜板／菱形板} ($\beta=15^{\circ},30^{\circ},45^{\circ}$) &
  [\ref{bathe1985}, \ref{onate1994}, \ref{krysl1995}, \ref{gu2001}, \ref{liew2003}, \ref{nguyenxuan2008}] &
  [\ref{bathe1985}, \ref{onate1994}, \ref{krysl1995}, \ref{gu2001}, \ref{liew2003}, \ref{nguyenxuan2008}] &
  [\ref{krysl1995}, \ref{gu2001}, \ref{nguyenxuan2008}] &
  --- &
  --- &
  --- \\
\midrule
\textbf{L型板／不規則板} (凹角) &
  [\ref{onate1994}, \ref{pian1969}, \ref{krysl1995}, \ref{liu1995}, \ref{liu2001}, \ref{gu2001}, \ref{liew2003}] &
  [\ref{onate1994}, \ref{pian1969}, \ref{krysl1995}, \ref{liu1995}, \ref{liu2001}, \ref{gu2001}, \ref{liew2003}] &
  --- &
  --- &
  --- &
  --- \\
\midrule
\textbf{圓環厚板} ($\frac{R_i}{R_o}=0.2,0.5$) &
  [\ref{bathe1985}, \ref{onate1994}, \ref{krysl1995}, \ref{liu1995}, \ref{liu2001}] &
  [\ref{bathe1985}, \ref{onate1994}, \ref{krysl1995}, \ref{liu1995}, \ref{liu2001}] &
  --- &
  --- &
  --- &
  --- \\
\bottomrule
\end{tabular}
\end{landscape}
\clearpage

% ─── 3. Final Output Items (new section) ─────────────────────────────────────
\section{有限元分析最終輸出項整理}
\label{sec:outputs}

本節彙整 Mindlin 板自由振動有限元分析中常見的**最終輸出項**，分為四個類別：
核心振動特性、力學響應與能量、數值評估、設計擴展。每個表格均列出輸出項的
常用符號、關鍵公式、解釋內容及其工程／科學意義。

\clearpage
\thispagestyle{empty}
\begin{landscape}
\centering
\footnotesize
\setlength{\tabcolsep}{4pt}
\renewcommand{\arraystretch}{1.15}

\subsection*{一、核心振動特性}

\begin{longtable}{p{2.2cm}p{2.2cm}p{5.5cm}p{4.5cm}p{4.5cm}}
\caption{核心振動特性輸出項} \label{tab:core_vib} \\
\toprule
\textbf{最終輸出項} & \textbf{常用符號} & \textbf{公式} & \textbf{核心解釋內容} & \textbf{工程／科學意義} \\
\midrule
\endfirsthead
\multicolumn{5}{c}{接續上頁} \\
\toprule
\textbf{最終輸出項} & \textbf{常用符號} & \textbf{公式} & \textbf{核心解釋內容} & \textbf{工程／科學意義} \\
\midrule
\endhead
\bottomrule
\endfoot
固有頻率 & $f$, $\omega$, $\Omega$ & $\omega=2\pi f$, $\Omega=\omega a^2\sqrt{\frac{\rho h}{D}}$, $D=\frac{Eh^3}{12(1-\nu^2)}$, $[K]\{\Phi\}=\omega^2[M]\{\Phi\}$ & 結構自然頻率，基頻最低；無量綱參數 $\Omega$ 消除尺寸、材料影響 & 衡量结构刚度与质量分布，特征值 $\lambda=\omega^2$反映了结构在特定模态下的刚度‑质量比。 \\
\addlinespace
模態振型 & $\Phi_i(x,y)$, $\{\Phi_i\}$ & $([K]-\omega_i^2[M])\{\Phi_i\}=0$, $\begin{Bmatrix} w\\\theta_x\\\theta_y\end{Bmatrix}=\sum_i\begin{Bmatrix} \Phi_{w,i}\\\Phi_{\theta_x,i}\\\Phi_{\theta_y,i}\end{Bmatrix}q_i(t)$ & 對應頻率的振動形態（節線、3D變形圖） & 验证空间变形模式的正确性，特征向量 ${\phi}$ 描述了结构在对应频率下的振动形态（挠度、转角分布）。通过可视化振型可以检查边界条件、对称性、非物理模态等問題 \\
\addlinespace
模態應力合力 & $M_x,M_y,M_{xy}$（彎矩）, $Q_x,Q_y$（剪力） & $\begin{Bmatrix}M\\Q\end{Bmatrix}=[D]\begin{Bmatrix}\kappa\\\gamma\end{Bmatrix}$, $\kappa$ 曲率，$\gamma$ 剪應變 & 模態對應的內力分佈（彎矩、剪力、扭矩） & 量化與预测裂纹尖端应力场与扩展方向，在裂纹尖端，应力理论上趋于无穷大（奇异性），因此无法直接用常规应力值进行评估，这在航空航天、压力容器、船舶结构中至关重要，用于制定检修周期或剩余寿命评估。 \\
\addlinespace
屈曲載荷 & $P_{cr}$, $\lambda_{cr}$ & $([K]+\lambda[K_G])\{U\}=0$, $P_{cr}=\lambda_{cr}P_{\text{ref}}$ & 面內壓力下失去穩定性的臨界值，考慮剪切變形 & 屈曲载荷与固有频率在数学上都归结为特征值问题。评估屈曲载荷的变化趋势有助于理解刚度退化与振动特性的关联，为动态稳定性分析提供基础。 \\
\addlinespace
誤差估計 & $e$, $\eta$, $E$ & $\eta=\sqrt{\frac{\|e\|_E^2}{\|\sigma\|_E^2}}$, $\|u-u_h\|_0=\bigl(\int_\Omega(u-u_h)^2d\Omega\bigr)^{1/2}$ & 數值解的精確度與收斂率，先驗／後驗估計 & 自適應細化、新單元驗證 \\
\addlinespace
應力強度因子 (SIF) & $K_I, K_{II}, K_{III}$ & $K=Y\sigma\sqrt{\pi a}$,  Mindlin 極限：$K_I^{\text{(RM)}}\to\frac{1+\nu}{3+\nu}K_I^{\text{(Kirchhoff)}}$ & 裂紋尖端應力場強度，透過 MVCCI 計算 & 裂紋擴展預測、剩餘壽命、無損檢測 \\
\end{longtable}
\end{landscape}
\normalsize

% ─── 4. Methodology Categories (originally 3) ────────────────────────────────
\section{依數值方法進行分類}
\label{sec:methodology}

所收集的文獻主要可歸納為以下七大數值方法群族：

\begin{enumerate}
    \itemsep0.3em
    \item \textbf{傳統位移場有限元素法} --- 採用等參數內插並搭配降階／選擇性積分
        以減輕剪力閉鎖效應~\ref{ahmad1970}, \ref{zienkiewicz1971}, 
        \ref{hughes1977}, \ref{pugh1978}, \ref{bathe1985}, \ref{belytschko1984}, 
        \ref{hinton1986}, \ref{onate1994}。
    \item \textbf{混合法與混合應力法} --- 透過對轉角與剪應變進行獨立內插，
        利用混合變分原理避開剪力閉鎖~\ref{pian1969}, \ref{bletzinger2000}。
    \item \textbf{無網格法 / 無元素 Galerkin（EFG）法} --- 基於移動最小二乘法
        或再生核函數進行近似，僅需節點資料即可建立離散化~\ref{belytschko1994},
        \ref{krysl1995}, \ref{liu1995}, \ref{chen1996}, \ref{wang2002},
        \ref{gu2001}, \ref{liu2001}。
    \item \textbf{無網格局部 Petrov--Galerkin（MLPG）法} --- 使用局部弱形式、
        真正不需要網格的方法~\ref{liew2003}, \ref{gu2001}。
    \item \textbf{等幾何分析法（IGA）} --- 基於 NURBS 的基底函數，整合 CAD 與
        分析流程，提供更高階的連續性（例如 $C^{1}$ 或 $C^{2}$），
        對厚板公式十分有利~\ref{hughes2005}, \ref{cottrell2006}, \ref{thai2012}。
    \item \textbf{光滑有限元素法（S\textit{FEM}）} --- 在節點、邊界或單元區域上
        進行應變平滑化，以軟化剛度並提高精度~\ref{liu2009}, \ref{nguyenxuan2008}。
    \item \textbf{擴展／富化有限元素法（XFEM）} --- 透過富化函數模擬板結構中
        的裂紋與不連續性~\ref{dolbow2000}。
\end{enumerate}

% ─── 5. Bibliographic References (originally 4) ──────────────────────────────
\section{完整參考文獻與 DOI 連結}
\label{sec:refs}

以下列出各篇文獻之完整書目資料與 DOI 超連結。編號對應表~\ref{tab:cases} 中所
使用的索引。

\begin{enumerate}

    \item \label{ahmad1970}
    S.~Ahmad, B.M.~Irons, and O.C.~Zienkiewicz.
    ``Analysis of thick and thin shell structures by curved finite elements.''
    \emph{International Journal for Numerical Methods in Engineering},
    vol.~2, no.~3, pp.~419--451, 1970.
    \newline \url{https://doi.org/10.1002/nme.1620020310}

    \item \label{zienkiewicz1971}
    O.C.~Zienkiewicz, R.L.~Taylor, and J.M.~Too.
    ``Reduced integration technique in general analysis of plates and shells.''
    \emph{International Journal for Numerical Methods in Engineering},
    vol.~3, no.~2, pp.~275--290, 1971.
    \newline \url{https://doi.org/10.1002/nme.1620030211}

    \item \label{hughes1977}
    T.J.R.~Hughes, R.L.~Taylor, and W.~Kanoknukulchai.
    ``A simple and efficient finite element for plate bending.''
    \emph{International Journal for Numerical Methods in Engineering},
    vol.~11, no.~10, pp.~1529--1543, 1977.
    \newline \url{https://doi.org/10.1002/nme.1620111005}

    \item \label{pugh1978}
    E.D.L.~Pugh, E.~Hinton, and O.C.~Zienkiewicz.
    ``A study of quadrilateral plate bending elements with reduced integration.''
    \emph{International Journal for Numerical Methods in Engineering},
    vol.~12, no.~7, pp.~1059--1079, 1978.
    \newline \url{https://doi.org/10.1002/nme.1620120702}

    \item \label{bathe1985}
    K.J.~Bathe and E.N.~Dvorkin.
    ``A four-node plate bending element based on Mindlin/Reissner plate theory 
    and a mixed interpolation.''
    \emph{International Journal for Numerical Methods in Engineering},
    vol.~21, no.~2, pp.~367--383, 1985.
    \newline \url{https://doi.org/10.1002/nme.1620210213}

    \item \label{belytschko1984}
    T.~Belytschko, J.I.~Lin, and C.S.~Tsay.
    ``Explicit algorithms for the nonlinear dynamics of shells.''
    \emph{Computer Methods in Applied Mechanics and Engineering},
    vol.~42, no.~2, pp.~225--251, 1984.
    \newline \url{https://doi.org/10.1016/0045-7825(84)90026-4}

    \item \label{hinton1986}
    E.~Hinton and H.C.~Huang.
    ``A family of quadrilateral Mindlin plate elements with substitute shear 
    strain fields.''
    \emph{Computers \& Structures}, vol.~23, no.~3, pp.~409--431, 1986.
    \newline \url{https://doi.org/10.1016/0045-7949(86)90233-5}

    \item \label{onate1994}
    E.~O\~nate, F.~Z\'arate, and F.~Flores.
    ``A basic thin/thick plate element.''
    \emph{International Journal for Numerical Methods in Engineering},
    vol.~37, no.~11, pp.~1949--1972, 1994.
    \newline \url{https://doi.org/10.1002/nme.1620371108}

    \item \label{bletzinger2000}
    K.U.~Bletzinger, M.~Bischoff, and E.~Ramm.
    ``A unified approach for shear-locking-free triangular and rectangular shell 
    elements.''
    \emph{Computers \& Structures}, vol.~75, no.~3, pp.~321--334, 2000.
    \newline \url{https://doi.org/10.1016/S0045-7949(99)00141-6}

    \item \label{pian1969}
    T.H.H.~Pian and P.~Tong.
    ``Basis of finite element methods for solid continua.''
    \emph{International Journal for Numerical Methods in Engineering},
    vol.~1, no.~1, pp.~3--28, 1969.
    \newline \url{https://doi.org/10.1002/nme.1620010103}

    \item \label{belytschko1994}
    T.~Belytschko, Y.Y.~Lu, and L.~Gu.
    ``Element-free Galerkin methods.''
    \emph{International Journal for Numerical Methods in Engineering},
    vol.~37, no.~2, pp.~229--256, 1994.
    \newline \url{https://doi.org/10.1002/nme.1620370205}

    \item \label{krysl1995}
    P.~Krysl and T.~Belytschko.
    ``Analysis of thin plates by the element-free Galerkin method.''
    \emph{Computational Mechanics}, vol.~17, no.~1--2, pp.~26--35, 1995.
    \newline \url{https://doi.org/10.1007/BF00356475}

    \item \label{liu1995}
    W.K.~Liu, S.~Jun, and Y.F.~Zhang.
    ``Reproducing kernel particle methods.''
    \emph{International Journal for Numerical Methods in Fluids},
    vol.~20, no.~8--9, pp.~1081--1106, 1995.
    \newline \url{https://doi.org/10.1002/fld.1650200824}

    \item \label{chen1996}
    J.S.~Chen, C.~Pan, C.T.~Wu, and W.K.~Liu.
    ``Reproducing kernel particle methods for large deformation analysis of 
    nonlinear structures.''
    \emph{Computer Methods in Applied Mechanics and Engineering},
    vol.~139, no.~1--4, pp.~195--227, 1996.
    \newline \url{https://doi.org/10.1016/S0045-7825(96)01083-3}

    \item \label{liu2001}
    G.R.~Liu and X.L.~Chen.
    ``A mesh-free method for static and free vibration analyses of thin plates 
    of complicated shape.''
    \emph{Journal of Sound and Vibration}, vol.~241, no.~5, pp.~839--855, 2001.
    \newline \url{https://doi.org/10.1006/jsvi.2000.3324}

    \item \label{wang2002}
    J.G.~Wang and G.R.~Liu.
    ``A point interpolation meshless method based on radial basis functions.''
    \emph{International Journal for Numerical Methods in Engineering},
    vol.~54, no.~11, pp.~1623--1648, 2002.
    \newline \url{https://doi.org/10.1002/nme.489}

    \item \label{gu2001}
    Y.T.~Gu and G.R.~Liu.
    ``A meshless local Petrov--Galerkin (MLPG) method for free and 
    forced vibration analyses for solids.''
    \emph{Computational Mechanics}, vol.~27, no.~3, pp.~188--198, 2001.
    \newline \url{https://doi.org/10.1007/s004660100232}

    \item \label{liew2003}
    K.M.~Liew, Y.Q.~Huang, and J.N.~Reddy.
    ``Vibration analysis of symmetrically laminated plates based on FSDT using 
    the moving least squares differential quadrature method.''
    \emph{Computer Methods in Applied Mechanics and Engineering},
    vol.~192, no.~19, pp.~2203--2222, 2003.
    \newline \url{https://doi.org/10.1016/S0045-7825(03)00238-X}

    \item \label{hughes2005}
    T.J.R.~Hughes, J.A.~Cottrell, and Y.~Bazilevs.
    ``Isogeometric analysis: CAD, finite elements, NURBS, exact geometry and 
    mesh refinement.''
    \emph{Computer Methods in Applied Mechanics and Engineering},
    vol.~194, no.~39--41, pp.~4135--4195, 2005.
    \newline \url{https://doi.org/10.1016/j.cma.2004.10.008}

    \item \label{cottrell2006}
    J.A.~Cottrell, A.~Reali, Y.~Bazilevs, and T.J.R.~Hughes.
    ``Studies of refinement and continuity in isogeometric structural analysis.''
    \emph{Computer Methods in Applied Mechanics and Engineering},
    vol.~196, no.~41--44, pp.~4160--4183, 2006.
    \newline \url{https://doi.org/10.1016/j.cma.2006.06.014}

    \item \label{thai2012}
    C.H.~Thai, H.~Nguyen-Xuan, N.~Nguyen-Thanh, T.H.~Le, T.~Nguyen-Thoi, 
    and T.~Rabczuk.
    ``Static, free vibration and buckling analysis of laminated composite 
    Reissner--Mindlin plates using NURBS-based isogeometric approach.''
    \emph{Computer Methods in Applied Mechanics and Engineering},
    vol.~247--248, pp.~10--33, 2012.
    \newline \url{https://doi.org/10.1016/j.cma.2012.07.018}

    \item \label{dolbow2000}
    J.~Dolbow, N.~Mo\"es, and T.~Belytschko.
    ``Modeling fracture in Mindlin--Reissner plates with the 
    extended finite element method.''
    \emph{International Journal of Solids and Structures},
    vol.~37, no.~48--50, pp.~7161--7183, 2000.
    \newline \url{https://doi.org/10.1016/S0020-7683(00)00194-5}

    \item \label{liu2009}
    G.R.~Liu, T.~Nguyen-Thoi, and K.Y.~Lam.
    ``An edge-based smoothed finite element method (ES-FEM) for static, free 
    and forced vibration analyses of solids.''
    \emph{Journal of Sound and Vibration}, vol.~320, no.~4--5, pp.~1100--1130, 2009.
    \newline \url{https://doi.org/10.1016/j.jsv.2008.08.027}

    \item \label{nguyenxuan2008}
    H.~Nguyen-Xuan, T.~Rabczuk, S.~Bordas, and J.F.~Debongnie.
    ``A smoothed finite element method for plate analysis.''
    \emph{Computer Methods in Applied Mechanics and Engineering},
    vol.~197, no.~13--16, pp.~1184--1203, 2008.
    \newline \url{https://doi.org/10.1016/j.cma.2007.09.022}

    \item \label{sukumar2001}
    N.~Sukumar, B.~Moran, A.Y.~Semenov, and V.V.~Belikov.
    ``Natural neighbour Galerkin methods.''
    \emph{International Journal for Numerical Methods in Engineering},
    vol.~50, no.~1, pp.~1--27, 2001.
    \newline \url{https://doi.org/10.1002/1097-0207(20010110)50:1<1::AID-NME14>3.0.CO;2-P}

\end{enumerate}

% ─── 6. Discussion (originally 5) ────────────────────────────────────────────
\section{討論}
\label{sec:discussion}

本回顧揭示了幾個重要的發展趨勢：

\begin{itemize}
    \item \textbf{剪力閉鎖（Shear locking）} 仍是低階 Mindlin 板元素的核心挑戰。
        目前已可透過 MITC（張量分量混合內插法）~\ref{bathe1985} 與 DSG
        （離散剪應變差法）~\ref{bletzinger2000} 等技術成功緩解。
    \item \textbf{無網格法}（EFG、RKPM、MLPG）可提供高精度及平滑的振態，
        但計算成本較高，且需在振動分析中加入穩定性處理。
    \item \textbf{等幾何分析法（IGA）} 近年來備受關注，因其高階連續性
        （$C^{1}$ 或 $C^{2}$）可自然滿足 Mindlin 理論的正則性要求，
        無需特殊處理~\ref{thai2012}。
    \item \textbf{光滑有限元素法}（ES-FEM、CS-FEM）在精度與計算效率之間
        取得了良好平衡，尤其適用於網格扭曲問題~\ref{nguyenxuan2008}。
    \item 大多數研究集中在 SSSS 或 CCCC 邊界條件的\textbf{方形板}，
        而斜板、L型板或不規則幾何的研究較少——這些是未來研究可填補的缺口。
\end{itemize}

% ─── 7. Conclusion (originally 6) ────────────────────────────────────────────
\section{結論}
本文針對 Mindlin 板自由振動分析之有限元素法進行了系統性的文獻回顧。
表~\ref{tab:cases} 所收錄的 25 篇參考文獻涵蓋七大數值方法家族、七種測試幾何
以及六種常見邊界條件組合。本彙整可為 \texttt{ApproxOperator.jl} 框架下
所開發的新方法提供方便的基準參考。

\end{document}
